%Two resources useful for abstract writing.
% Guidance of how to write an abstract/summary provided by Nature: https://cbs.umn.edu/sites/cbs.umn.edu/files/public/downloads/Annotated_Nature_abstract.pdf %https://writingcenter.gmu.edu/guides/writing-an-abstract
\chapter*{\center \Large  Abstract}
% The advancement of autonomous Unmanned Aerial Vehicles (UAVs) requires control strategies that enables high speed maneuvering with rigorous stability guarantees in the presence of environmental disturbances. This research explores the development and comparative evaluation of different control paradigms applied to the Crazyflie 2.1 nano-quadrotor platform. The study initiates with the derivation of a six degree of freedom (6-DOF) dynamical model to capture the nonlinearities and coupling effects inherent in the system's translational and rotational dynamics. A cascaded Proportional-Integral-Derivative (PID) controller is first established as a linear baseline, followed by the implementation of three nonlinear strategies: Conventional Sliding Mode Control (SMC), Super-Twisting Sliding Mode Control (STSMC), and Non-singular Super-Twisting Terminal Sliding Mode Control (NSTSMC). Each progression addresses specific limitations of the previous architecture, specifically focusing on the rejection of aerodynamic disturbances, the mitigation of high-frequency chattering, and the achievement of finite-time convergence while avoiding mathematical singularities. The control laws are first validated through numerical simulations in MATLAB, focusing on steady-state accuracy and transient response. Subsequently, the research transitions to a high-fidelity validation phase using the Robot Operating System (ROS) and the Gazebo simulation environment, employing the Crazyswarm2 framework and the ros-gz-crazyflie bridge. Results indicate that the NSTSMC architecture provides the most robust performance, demonstrating superior trajectory tracking and disturbance rejection capabilities under varying load conditions and external wind models.

The advancement of autonomous Unmanned Aerial Vehicles (UAVs) pose challenging nonlinear control problems due to underactuation and disturbances (e.g. wind, model uncertainty). This work derives the complete 6-degree-of-freedom quadrotor dynamics from first principles (Newton–Euler formulation) to capture its coupling and gravity effects. A hierarchical control strategy is investigated. First, a standard cascaded PID controller is implemented (inner loop for attitude stabilization, outer loop for position) as a baseline. Next, a conventional sliding-mode controller (SMC) is designed: attitude errors define sliding surfaces and SMC laws drive these surfaces to zero, guaranteeing robustness to parameter variation. To overcome the high-frequency chattering of first-order SMC, we introduce a super-twisting sliding-mode controller (STSMC), which uses a continuous equivalent control to suppress chattering while retaining disturbance rejection. Finally, a non-singular super-twisting terminal SMC (NSTSMC) that ensures finite-time convergence of errors without singularities is designed. Each control law is derived with Lyapunov based stability proofs.

The controllers are tested in simulation and validated experimentally. MATLAB simulations run aggressive trajectory following tasks under payload and wind disturbances; performance is measured by various metrics. The models and controllers are then tested on a Crazyflie 2.1 quadrotor model using ROS and Gazebo to emulate hardware. All designs maintain stability; however, the advanced SMCs consistently outperform the PID baseline. In particular, STSMC significantly reduces chattering and improves disturbance rejection, while NSTSMC achieves the fastest convergence and tightest tracking despite uncertainties.
% For example, NSTSMC yields qualitatively shorter settling time and smaller steady-state error (exact values depend on tuning). 

This thesis contributes to a detailed quadrotor dynamic model, a family of sliding-mode controllers with increasing sophistication, and  systematic validation in both simulation and hardware-in-loop. The comparative analysis shows the proposed NSTSMC controller provides the best robustness and performance in uncertain conditions. Quantitative performance improvements are summarized qualitatively in the results. These findings guide the design of robust controllers for agile UAVs in real environments. 